
\documentclass{article}
\usepackage{colortbl}
\usepackage{makecell}
\usepackage{multirow}
\usepackage{supertabular}

\begin{document}

\newcounter{utterance}

\centering \large Interaction Transcript for game `cladder', experiment `full\_v1.5\_default', episode 6357 with qwen3.5{-}27b.
\vspace{24pt}

{ \footnotesize  \setcounter{utterance}{1}
\setlength{\tabcolsep}{0pt}
\begin{supertabular}{c@{$\;$}|p{.15\linewidth}@{}p{.15\linewidth}p{.15\linewidth}p{.15\linewidth}p{.15\linewidth}p{.15\linewidth}}
    \# & \multicolumn{2}{c}{Player} && \multicolumn{2}{c}{Game Master} \\
    \hline

    \theutterance \stepcounter{utterance}  
    & & & \multicolumn{4}{p{0.6\linewidth}}{
        \cellcolor[rgb]{0.9,0.9,0.9}{
            \makecell[{{p{\linewidth}}}]{
                \texttt{\tiny{[P1$\langle$GM]}}
                \texttt{You are an expert in causal inference. The following question is not a typical commonsense query, but rather a meticulously designed question created by a professor specializing in causal inference, intended to assess the students' mastery of the course content.} \\
\texttt{Imagine a self{-}contained, hypothetical world with only the following conditions, and without any unmentioned factors or causal relationships: Rixq has a direct effect on zuph and xevu. Swoy has a direct effect on zuph and xevu. Zuph has a direct effect on xevu. Swoy is unobserved. The overall probability of rixq is 70\%. The probability of not rixq and xevu is 6\%. The probability of rixq and xevu is 10\%. Is the chance of xevu larger when observing rixq?} \\
\texttt{Start your answer with 'Yes' or 'No', followed by additional reasoning or evidence to support your explanation.} \\
            }
        }
    }
     \\ \\

    \theutterance \stepcounter{utterance}  
    & \multicolumn{4}{p{0.6\linewidth}}{
        \cellcolor[rgb]{0.561,0.737,0.561}{
            \makecell[{{p{\linewidth}}}]{
                \texttt{\tiny{[P1$\rangle$GM]}}
                \texttt{No} \\
\\ 
\texttt{To determine whether the chance of \$xevu\$ is larger when observing \$rixq\$, we must calculate the conditional probabilities \$P(xevu \textbackslash{}mid rixq)\$ and \$P(xevu \textbackslash{}mid \textbackslash{}neg rixq)\$ based strictly on the provided numerical data.} \\
\\ 
\texttt{**1. Extract Given Probabilities:**} \\
\texttt{*   \$P(rixq) = 0.70\$} \\
\texttt{*   \$P(\textbackslash{}neg rixq) = 1 {-} 0.70 = 0.30\$} \\
\texttt{*   \$P(\textbackslash{}neg rixq \textbackslash{}cap xevu) = 0.06\$} \\
\texttt{*   \$P(rixq \textbackslash{}cap xevu) = 0.10\$} \\
\\ 
\texttt{**2. Calculate Conditional Probabilities:**} \\
\\ 
\texttt{*   **Probability of \$xevu\$ given \$rixq\$:**} \\
\texttt{    \$\$P(xevu \textbackslash{}mid rixq) = \textbackslash{}frac\{P(rixq \textbackslash{}cap xevu)\}\{P(rixq)\} = \textbackslash{}frac\{0.10\}\{0.70\} \textbackslash{}approx 0.143\$\$} \\
\\ 
\texttt{*   **Probability of \$xevu\$ given \$\textbackslash{}neg rixq\$:**} \\
\texttt{    \$\$P(xevu \textbackslash{}mid \textbackslash{}neg rixq) = \textbackslash{}frac\{P(\textbackslash{}neg rixq \textbackslash{}cap xevu)\}\{P(\textbackslash{}neg rixq)\} = \textbackslash{}frac\{0.06\}\{0.30\} = 0.20\$\$} \\
\\ 
\texttt{**3. Comparison and Causal Context:**} \\
\texttt{Comparing the two values:} \\
\texttt{\$\$0.143 < 0.20\$\$} \\
\texttt{Therefore, \$P(xevu \textbackslash{}mid rixq) < P(xevu \textbackslash{}mid \textbackslash{}neg rixq)\$. The chance of \$xevu\$ is actually **smaller** when observing \$rixq\$.} \\
\\ 
\texttt{**Causal Interpretation:**} \\
\texttt{Although the problem states that \$rixq\$ has a *direct effect* on \$xevu\$, the observed association is negative. This is a classic example of **confounding bias** (specifically, a collider or confounder structure depending on the exact DAG, but here \$Swoy\$ acts as a confounder).} \\
\texttt{*   \$Swoy\$ affects both \$rixq\$ and \$xevu\$.} \\
\texttt{*   \$Swoy\$ is **unobserved**.} \\
\texttt{*   Because we cannot control for \$Swoy\$, the marginal association between \$rixq\$ and \$xevu\$ is distorted. The data shows that when \$rixq\$ is present, \$xevu\$ is less likely, likely because the presence of \$rixq\$ is correlated with a state of \$Swoy\$ that suppresses \$xevu\$, or the specific parameterization of the unobserved confounder creates this negative correlation despite the direct positive causal link.} \\
\\ 
\texttt{The question asks about the "chance... when observing," which refers to the conditional probability derived from the data, not the counterfactual causal effect. Based on the data provided, the probability decreases.} \\
            }
        }
    }
    & & \\ \\

    \theutterance \stepcounter{utterance}  
    & & & \multicolumn{2}{p{0.3\linewidth}}{
        \cellcolor[rgb]{0.9,0.9,0.9}{
            \makecell[{{p{\linewidth}}}]{
                \texttt{\tiny{[GM$|$GM]}}
                \texttt{No} \\
            }
        }
    }
    & & \\ \\

    \theutterance \stepcounter{utterance}  
    & & & \multicolumn{2}{p{0.3\linewidth}}{
        \cellcolor[rgb]{0.9,0.9,0.9}{
            \makecell[{{p{\linewidth}}}]{
                \texttt{\tiny{[GM$|$GM]}}
                \texttt{no} \\
            }
        }
    }
    & & \\ \\

    \theutterance \stepcounter{utterance}  
    & & & \multicolumn{2}{p{0.3\linewidth}}{
        \cellcolor[rgb]{0.9,0.9,0.9}{
            \makecell[{{p{\linewidth}}}]{
                \texttt{\tiny{[GM$|$GM]}}
                \texttt{game\_result = WIN} \\
            }
        }
    }
    & & \\ \\

\end{supertabular}
}

\end{document}
